\chapter{COVID-19}

\section*{Definition and Overview}
COVID-19 is an infectious disease caused by the SARS-CoV-2 virus, a coronavirus that emerged in late 2019 and caused a global pandemic. The virus spreads mainly through respiratory droplets and aerosols produced when an infected person coughs, sneezes, talks, or breathes. COVID-19 causes a wide range of illness, from no symptoms at all to mild cold-like symptoms, severe pneumonia, and life-threatening complications such as respiratory failure and blood clots. Vaccines, antiviral treatments, and improved immunity have greatly reduced the severity of the pandemic, but COVID-19 remains an ongoing public health concern in Australia and worldwide. This chapter covers transmission, symptoms, testing, treatment, and prevention.

\section*{Epidemiology}
COVID-19 has infected hundreds of millions of people worldwide and caused millions of deaths, although the toll has fallen sharply since the introduction of vaccines and treatments. In Australia, the pandemic progressed through distinct waves driven by new variants, from the original strain to Delta and then Omicron and its subvariants. Most of the Australian population has now been infected, vaccinated, or both. The disease is generally milder in vaccinated people and in children, and is most severe in older adults, people with chronic conditions, and those who are immunocompromised. Long COVID, a post-infection syndrome, affects a minority of people after recovery.

\section*{Aetiology and Pathophysiology}
\begin{itemize}
    \item \textbf{Virus:} SARS-CoV-2 is a single-stranded RNA coronavirus; its spike protein binds to ACE2 receptors on human cells
    \item \textbf{Transmission:} mainly by respiratory droplets and aerosols; also via contaminated surfaces, though this is less important
    \item \textbf{Incubation:} symptoms usually appear 2--14 days after exposure, most often around 4--6 days
    \item \textbf{Replication:} the virus multiplies in the airways, triggering an immune response and inflammation
    \item \textbf{Severe disease:} in some people the immune response becomes dysregulated, causing lung damage, blood clotting, and multi-organ injury
    \item \textbf{Immunity:} vaccination and prior infection provide protection against severe illness, though reinfection occurs
\end{itemize}

\section*{Clinical Features}
\begin{itemize}
    \item Fever, cough, sore throat, and runny nose
    \item Fatigue, headache, and muscle aches
    \item Loss of smell or taste, which was characteristic of early variants
    \item Shortness of breath, which indicates more serious lung involvement
    \item Nausea, vomiting, and diarrhoea in some people
    \item Many people, especially children and vaccinated adults, have no symptoms or a mild illness resembling a cold
    \item Severe features: breathlessness, chest pain, confusion, and low oxygen levels
    \item Long COVID: persistent fatigue, brain fog, and other symptoms lasting months after the acute infection
\end{itemize}

\section*{Differential Diagnosis}
\begin{itemize}
    \item Influenza and other respiratory viruses, which cause similar symptoms
    \item The common cold and other coronaviruses
    \item Bacterial pneumonia, which may follow or coexist with viral infection
    \item Respiratory syncytial virus (RSV) and other childhood respiratory infections
    \item Allergic rhinitis and asthma exacerbations
    \item Other causes of fever, including urinary or other infections
\end{itemize}

\section*{When to Seek Medical Care}
People with symptoms should test for COVID-19 and stay home while unwell. Those at high risk, including older adults, pregnant women, and people with chronic conditions, should contact their doctor about early antiviral treatment, which works best when started promptly. Seek medical review if symptoms are severe, worsening, or not improving.

\textbf{Call 000 (or local emergency services) for:}
\begin{itemize}
    \item Difficulty breathing, shortness of breath at rest, or chest pain
    \item Confusion, drowsiness, or difficulty waking
    \item Bluish lips or face, or a sudden drop in oxygen levels
    \item Severe dehydration with inability to keep down fluids
    \item Seizures or collapse
\end{itemize}

\section*{Diagnosis and Investigations}
\begin{itemize}
    \item \textbf{Rapid antigen tests (RATs):} home tests that detect the virus; useful for diagnosis during the infectious period
    \item \textbf{PCR tests:} more sensitive laboratory tests, used to confirm infection and in healthcare settings
    \item \textbf{Pulse oximetry:} measures blood oxygen levels and helps identify severe disease at home
    \item \textbf{Chest X-ray or CT:} shows lung changes in people with pneumonia
    \item \textbf{Blood tests:} inflammatory markers, organ function, and clotting studies in severe disease
    \item \textbf{Assessment of risk:} age, vaccination status, and medical conditions guide treatment decisions
\end{itemize}

\section*{Management}
\begin{itemize}
    \item \textbf{Self-care at home:} rest, fluids, and paracetamol for fever and pain
    \item \textbf{Antiviral medicines:} oral treatments such as nirmatrelvir-ritonavir (Paxlovid) and molnupiravir are available for people at high risk, most effective when started within five days of symptoms
    \item \textbf{Isolation:} staying home while infectious reduces spread to others
    \item \textbf{Hospital care:} oxygen, corticosteroids, and other treatments for severe illness
    \item \textbf{Management of long COVID:} rehabilitation, symptom management, and specialist clinics
    \item \textbf{Monitoring:} watching for warning signs and seeking care if they develop
    \item \textbf{Immunisation:} staying up to date with COVID-19 vaccines reduces the risk of severe disease
\end{itemize}

\section*{Complications}
\begin{itemize}
    \item Pneumonia and acute respiratory distress syndrome (ARDS)
    \item Blood clots affecting the lungs, legs, or brain
    \item Myocarditis and other heart problems
    \item Kidney injury, sepsis, and multi-organ failure
    \item Long COVID, with persistent fatigue, brain fog, and other symptoms
    \item Increased risk of complications in pregnancy
    \item Death, mainly in older and high-risk people
\end{itemize}

\section*{Prognosis and Outcomes}
The outlook for COVID-19 has improved dramatically since the start of the pandemic. The vast majority of people, especially those who are vaccinated and healthy, recover fully within one to two weeks. Severe disease is now uncommon among vaccinated people but still occurs in older adults and those with significant risk factors. Most people hospitalised with COVID-19 recover, though recovery from severe illness can take weeks. Long COVID resolves in most affected people over months, but a minority have persisting symptoms. Staying up to date with vaccination remains the most effective way to reduce the risk of serious illness.

\section*{Prevention and Self-Care}
\begin{itemize}
    \item Stay up to date with recommended COVID-19 vaccinations, including boosters
    \item Stay home when unwell and test to know your status
    \item Practise good hand hygiene and cover coughs and sneezes
    \item Improve indoor ventilation, and consider masks in crowded indoor settings during periods of high transmission
    \item Discuss early antiviral treatment with a doctor if you are at high risk
    \item Maintain general health with good nutrition, sleep, and activity
    \item Follow current Australian health department advice on isolation and testing
\end{itemize}

\section*{Sources and Further Reading}
The following reputable sources provide further information for healthcare students and patients seeking reliable, current advice about COVID-19.

\begin{itemize}
    \item Healthdirect Australia. \textit{COVID-19: symptoms, testing, and treatment.} https://www.healthdirect.gov.au/covid-19
    \item World Health Organization. \textit{Coronavirus disease (COVID-19) fact sheet.} https://www.who.int/
    \item Australian Government Department of Health. \textit{COVID-19 vaccines and health information.} https://www.health.gov.au/
    \item NHS. \textit{COVID-19: symptoms and treatment.} https://www.nhs.uk/conditions/covid-19/
    \item Mayo Clinic. \textit{COVID-19: diagnosis and management.} https://www.mayoclinic.org/
    \item NSW Health. \textit{COVID-19 information for the public.} https://www.health.nsw.gov.au/
\end{itemize}
